%% toc: P1 4 4.1,4.2,4.3,4.4
\section{Injectivity of the numerical coding of strings}
\label{sec:p01-w02-p01-012-string-coding}

The countability lemma of \cref{sec:strings-01} shows that \(A^*\) is
countable whenever \(A\) is, by coding each string by a natural number. Its
proof makes two claims in one sentence each: that every stage map \(e_n\) is
injective, and that tagging the stage code with the length gives an injection
\(A^*\to\N\) \cite[ch.~``Finite Strings, Induction, and Recursion'',
sec.~``Finite sequences and strings'', lemma ``Countability of finite
strings'']{FrozenFormalLogicStringCoding}. This section writes out every step
of both claims. The statement, the coding maps and their notation are those
of the lemma; the step-by-step arguments, the decoding procedure, the worked
example and the failure case are an editorial expansion of its proof, not text
attributed to the source.

\emph{Inputs.} A set \(A\), an injection \(e:A\to\N\), and the numerical
pairing map \(\operatorname{pair}_{\N}\) fixed in \cref{sec:strings-01}.
\emph{Conclusions.} Each \(e_n\) is an injection \(A^n\to\N\)
(\cref{thm:p01-w02-p01-012-stage-injective}); every string has exactly one
length, and the length-tagged map is an injection \(A^*\to\N\)
(\cref{thm:p01-w02-p01-012-tagged-injective}); hence the countability lemma
(\cref{thm:p01-w02-p01-012-countable}); and a decoding procedure recovers a
string from its code (\cref{def:p01-w02-p01-012-decoding},
\cref{thm:p01-w02-p01-012-decoding-correct}).

\subsection{Definitions and assumptions}
\label{sec:p01-w02-p01-012-assumptions}

A finite sequence in \(A\) of length \(n\) is a function
\(w:\{0,\ldots,n-1\}\to A\); its length is \(\operatorname{len}(w)=n\);
\(\epsilon\) is the unique sequence of length \(0\); \(A^n\) is the set of
sequences in \(A\) of length \(n\), and \(A^*=\bigcup_{n\in\N}A^n\)
\cite[ch.~``Finite Strings, Induction, and Recursion'', definitions ``Finite
sequence'' and ``String'']{FrozenFormalLogicStringCoding}. A set is
\emph{countable} when there is an injection from it into \(\N\); finite sets
count as countable (\cref{sec:set-prerequisites-05}).

The pairing map is
\[
  \operatorname{pair}_{\N}(a,b)=\frac{(a+b)(a+b+1)}2+b
  \qquad(a,b\in\N),
\]
and \cref{sec:strings-01} proves that it is a bijection
\(\N\times\N\to\N\), using the triangular numbers \(t_s=s(s+1)/2\): for each
\(z\in\N\) there is exactly one \(s\) with \(t_s\leq z<t_{s+1}\), and then
\(b=z-t_s\), \(a=s-b\). The proofs below use only the injectivity half of
that fact, which they write out in the form needed.

The source proof writes \(w\mathbin{\upharpoonright}n\) for
\(w\in A^{n+1}\). This is restriction of the function \(w\) to the set
\(\{0,\ldots,n-1\}\): the sequence of length \(n\) that agrees with \(w\) at
every argument below \(n\). Two finite sequences are equal exactly when they
have the same length and the same value at every argument, because a function
is fixed by its domain and its values.

\begin{hbdefinition}{def:p01-w02-p01-012-coding-maps}{Stage codes and the length-tagged code}
Let \(e:A\to\N\) be an injection. The \emph{stage codes}
\(e_n:A^n\to\N\) are defined by recursion on \(n\):
\[
  e_0(\epsilon)=0,
  \qquad
  e_{n+1}(w)=\operatorname{pair}_{\N}\bigl(e_n(w\mathbin{\upharpoonright}n),
    e(w(n))\bigr)\quad(w\in A^{n+1}).
\]
The \emph{length-tagged code} sends \(w\in A^n\) to
\(\operatorname{pair}_{\N}(n,e_n(w))\).
\end{hbdefinition}

The recursion is well formed. For \(w\in A^{n+1}\) the restriction
\(w\mathbin{\upharpoonright}n\) has domain \(\{0,\ldots,n-1\}\) and values in
\(A\), so it lies in \(A^n\), where \(e_n\) is already defined; and
\(n\) belongs to the domain \(\{0,\ldots,n\}\) of \(w\), so \(w(n)\in A\)
and \(e(w(n))\) is defined. Only the one injection \(e\) is fixed; no choice
is made at any later stage.

\subsection{Results}
\label{sec:p01-w02-p01-012-results}

\begin{hblemma}{thm:p01-w02-p01-012-stage-injective}{proved=annex:p01-w02-p01-012-stage-injective}{Every stage code is injective}
Let \(e:A\to\N\) be an injection. For every \(n\in\N\), the stage code
\(e_n:A^n\to\N\) is injective.
\end{hblemma}

The proof is an induction on \(n\). At \(n=0\) the set \(A^0\) has the single
element \(\epsilon\). For the step, an equation
\(e_{n+1}(w)=e_{n+1}(w')\) is split by the injectivity of
\(\operatorname{pair}_{\N}\) into
\(e_n(w\mathbin{\upharpoonright}n)=e_n(w'\mathbin{\upharpoonright}n)\) and
\(e(w(n))=e(w'(n))\); the induction hypothesis and the injectivity of \(e\)
then give agreement below \(n\) and at \(n\).

\begin{hblemma}{thm:p01-w02-p01-012-tagged-injective}{proved=annex:p01-w02-p01-012-tagged-injective}{The length-tagged code is an injection}
Let \(e:A\to\N\) be an injection.
\begin{enumerate}
  \item If \(m\neq n\), then \(A^m\cap A^n=\emptyset\). Hence each
        \(w\in A^*\) lies in \(A^n\) for exactly one \(n\), namely
        \(n=\operatorname{len}(w)\), and the length-tagged code is a function
        on \(A^*\).
  \item If \(w\in A^n\), \(w'\in A^{n'}\) and
        \(\operatorname{pair}_{\N}(n,e_n(w))
          =\operatorname{pair}_{\N}(n',e_{n'}(w'))\), then \(n=n'\) and
        \(w=w'\). Thus the length-tagged code is an injection
        \(A^*\to\N\).
\end{enumerate}
\end{hblemma}

\begin{hbcorollary}{thm:p01-w02-p01-012-countable}{proved=annex:p01-w02-p01-012-countable}{Countability of finite strings}
If \(A\) is countable, then \(A^*\) is countable.
\end{hbcorollary}

This is the source lemma itself
\cite[ch.~``Finite Strings, Induction, and Recursion'', lemma
``Countability of finite strings'']{FrozenFormalLogicStringCoding}, here with
its proof assembled from
\cref{thm:p01-w02-p01-012-stage-injective,thm:p01-w02-p01-012-tagged-injective}.

\subsection{Decoding procedure}
\label{sec:p01-w02-p01-012-decoding-procedure}

The injectivity proofs can be run as a procedure that reads a string back off
its code. For \(z\in\N\) write \(\operatorname{unpair}(z)=(a,b)\) for the
inverse values of \cref{sec:strings-01}: \(s\) is the unique index with
\(t_s\leq z<t_{s+1}\), \(b=z-t_s\) and \(a=s-b\). For a letter \(x\in A\),
write \(\langle x\rangle\) for the string of length \(1\) whose value at
\(0\) is \(x\); concatenation \(uv\) is that of \cref{sec:strings-01}.

\begin{hbmethod}{def:p01-w02-p01-012-decoding}{Decoding a length-tagged code}
Input: an injection \(e:A\to\N\) and a number \(z\in\N\).
\begin{enumerate}
  \item Compute \((n,m)=\operatorname{unpair}(z)\).
  \item Compute \(d_n(m)\), where the partial maps \(d_k\) from \(\N\) to
        \(A^k\) are given by recursion on \(k\):
        \(d_0(0)=\epsilon\), and \(d_0\) is undefined at every other number;
        \(d_{k+1}(y)\) is computed by putting
        \((a,b)=\operatorname{unpair}(y)\) and returning
        \(d_k(a)\langle x\rangle\) for the unique \(x\in A\) with
        \(e(x)=b\).
  \item Output \(d_n(m)\).
\end{enumerate}
Side conditions. The letter \(x\) in step 2 is unique because \(e\) is
injective. If \(b\) is not a value of \(e\), or if \(d_k(a)\) is undefined,
then \(d_{k+1}(y)\) is undefined, and the procedure stops without output: the
number \(z\) is then not a length-tagged code. Running the procedure as an
algorithm also needs a way to compute \(e\) and to recognise and invert its
values; the existence claims of this section need only the function \(e\).
\end{hbmethod}

\begin{hblemma}{thm:p01-w02-p01-012-decoding-correct}{proved=annex:p01-w02-p01-012-decoding-correct}{Decoding recovers the string}
Let \(e:A\to\N\) be an injection. For every \(n\in\N\) and every
\(w\in A^n\), \(d_n(e_n(w))\) is defined and equals \(w\). Consequently the
procedure of \cref{def:p01-w02-p01-012-decoding}, applied to
\(z=\operatorname{pair}_{\N}(n,e_n(w))\), outputs \(w\).
\end{hblemma}

A map with a left inverse is injective, so
\cref{thm:p01-w02-p01-012-decoding-correct} gives a second route to the
injectivity of the stage codes and of the length-tagged code. The first route
(\cref{thm:p01-w02-p01-012-stage-injective,thm:p01-w02-p01-012-tagged-injective})
is the one the source proof indicates.

\subsection{Worked example and failure case}
\label{sec:p01-w02-p01-012-example}

\begin{hbexample}{def:p01-w02-p01-012-two-letters}{Codes over a two-letter alphabet}
Let \(A=\{\mathtt a,\mathtt b\}\) and \(e(\mathtt a)=0\),
\(e(\mathtt b)=1\). Write strings by listing their values, so
\(\mathtt{ab}\) is the sequence with value \(\mathtt a\) at \(0\) and
\(\mathtt b\) at \(1\). The pairing values used, computed from
\((a+b)(a+b+1)/2+b\), are
\[
\begin{aligned}
  &\operatorname{pair}_{\N}(0,0)=0,\quad
  \operatorname{pair}_{\N}(0,1)=2,\quad
  \operatorname{pair}_{\N}(1,0)=1,\\
  &\operatorname{pair}_{\N}(2,0)=3,\quad
  \operatorname{pair}_{\N}(2,1)=7,\quad
  \operatorname{pair}_{\N}(1,2)=8,\\
  &\operatorname{pair}_{\N}(2,2)=12,\quad
  \operatorname{pair}_{\N}(2,3)=18,\quad
  \operatorname{pair}_{\N}(2,7)=52.
\end{aligned}
\]
\begin{center}
\begin{tabular}{lll}
  \(w\) & \(e_n(w)\), \(n=\operatorname{len}(w)\)
        & \(\operatorname{pair}_{\N}(n,e_n(w))\)\\\hline
  \(\epsilon\) & \(0\) & \(0\)\\
  \(\mathtt a\) & \(\operatorname{pair}_{\N}(0,0)=0\) & \(1\)\\
  \(\mathtt b\) & \(\operatorname{pair}_{\N}(0,1)=2\) & \(8\)\\
  \(\mathtt{aa}\) & \(\operatorname{pair}_{\N}(0,0)=0\) & \(3\)\\
  \(\mathtt{ab}\) & \(\operatorname{pair}_{\N}(0,1)=2\) & \(12\)\\
  \(\mathtt{ba}\) & \(\operatorname{pair}_{\N}(2,0)=3\) & \(18\)\\
  \(\mathtt{bb}\) & \(\operatorname{pair}_{\N}(2,1)=7\) & \(52\)
\end{tabular}
\end{center}
For example
\(e_2(\mathtt{ba})=\operatorname{pair}_{\N}(e_1(\mathtt b),e(\mathtt a))
=\operatorname{pair}_{\N}(2,0)=3\). Within each length the stage codes are
distinct, as \cref{thm:p01-w02-p01-012-stage-injective} requires, and all
seven tagged codes are distinct.

Decoding \(z=12\): \(t_4=10\leq12<15=t_5\), so \(s=4\), \(b=2\), \(a=2\), and
\((n,m)=(2,2)\). Next \(\operatorname{unpair}(2)=(0,1)\), since
\(t_1=1\leq2<3=t_2\); the letter with code \(1\) is \(\mathtt b\). Then
\(\operatorname{unpair}(0)=(0,0)\); the letter with code \(0\) is
\(\mathtt a\), and \(d_0(0)=\epsilon\). So
\(d_2(2)=\epsilon\langle\mathtt a\rangle\langle\mathtt b\rangle=\mathtt{ab}\).

Decoding \(z=2\): \(\operatorname{unpair}(2)=(0,1)\), so \(n=0\) and
\(m=1\), and \(d_0(1)\) is undefined. The number \(2\) is not a tagged code:
the length-tagged code is injective but not surjective, which is all the
definition of countability asks for.
\end{hbexample}

\begin{hbwarning}{def:p01-w02-p01-012-untagged}{Dropping the length tag or the injectivity of the letter code}
\emph{Without the length tag.} The stage codes of different lengths can
coincide, so the union of the maps \(e_n\) need not be injective on \(A^*\).
Whenever some letter \(x\) has \(e(x)=0\),
\(e_1(\langle x\rangle)=\operatorname{pair}_{\N}(0,0)=0=e_0(\epsilon)\). In
\cref{def:p01-w02-p01-012-two-letters},
\(e_0(\epsilon)=e_1(\mathtt a)=e_2(\mathtt{aa})=0\) and
\(e_1(\mathtt b)=e_2(\mathtt{ab})=2\). The tag \(n\) is what separates the
lengths in part 2 of \cref{thm:p01-w02-p01-012-tagged-injective}.

\emph{Without injectivity of \(e\).} If \(e(x)=e(y)\) for letters \(x\neq y\),
then \(e_1(\langle x\rangle)=\operatorname{pair}_{\N}(0,e(x))
=\operatorname{pair}_{\N}(0,e(y))=e_1(\langle y\rangle)\), so \(e_1\) is
already not injective. The hypothesis that \(A\) is countable is used exactly
to supply an injective \(e\).

\emph{Limitation.} The construction gives an injection, not a bijection, as
the code \(2\) above shows. For a nonempty countable \(A\) the set \(A^*\) is
infinite (it contains a string of every length), and a bijection
\(\N\to A^*\) requires a further argument, such as the
Cantor--Schr\"oder--Bernstein theorem of \cref{sec:set-prerequisites-05}
together with an injection \(\N\to A^*\).
\end{hbwarning}

\subsection{Relationships and sources}
\label{sec:p01-w02-p01-012-relationships}

The corollary of \cref{sec:strings-01} on the countability of formulas over a
countable alphabet applies \cref{thm:p01-w02-p01-012-countable} to a subset of
\(A^*\). The same pairing map codes finite tuples of numbers for machines in
\cref{sec:turing-machines-01}, and the arithmetic coding of syntax builds its
constructor \(C(a,b)=1+\operatorname{pair}_{\N}(a,b)\) from it in
\cref{sec:arithmetic-coding-04}; both rely on the injectivity argument
written here in its simplest setting.

Sources. The statement, the definitions of finite sequence and string, the
pairing map with its bijectivity argument, and the recursion defining the
stage codes are from the held source \cite[ch.~``Finite Strings, Induction,
and Recursion'', sec.~``Finite sequences and strings'']{FrozenFormalLogicStringCoding};
the definition of countable set is from \cite[ch.~``Sets, Relations, Functions, and Cardinality'',
sec.~``Finite, countable, and uncountable sets'', definition ``Countable
set'']{FrozenFormalLogicStringCoding}. The written-out induction, the
disjointness of the sets \(A^n\), the decoding procedure, the example and the
failure case are editorial.
